\def\chaptercode{AL03}    % Code du chapitre
\def\chaptername{Diviser pour régner}        % Nom du chapitre

\titre{}

\medskip

\section{L'algorithme de dichotomie}

\subsection{La version impérative}

\begin{exercice}{Dichotomie impérative}{}
\medskip

Ecrire la fonction \verb+recherche_dicho_imperative(tab, val)+ qui renvoie \verb+True+ si \verb+val+ est dans \verb+tab+ et \verb+False+ sinon.

\begin{minipage}{0.49\linewidth}
    \begin{CodePythontexAlt}[Largeur=1\linewidth,Centre,Lignes=true,PremLigne=1]{}
        def recherche_dicho_imperative(tab, val):
            pass
        
        tab = [1, 5, 7, 9, 12, 13]
        assert recherche_dichotomique(tab, 12) == True
        assert recherche_dichotomique(tab, 17) == False
                \end{CodePythontexAlt}
\end{minipage}\hfill
\begin{minipage}{0.49\linewidth}
    À chaque tour de la boucle while, la taille de la liste est divisée par 2. Ceci confère à cet algorithme une complexité logarithmique (bien meilleure qu'une complexité linéaire).

\end{minipage}


\end{exercice}

%%%%%%%%%%%%%%%%%%%%%%%%%%%%%%%%%%%%%%%%%%%%%%%%%%%%%%%%%%%%%%%%%%%%%%%%%%%%%%%%%%%%%%%%%%%%%%%%%


\subsection{Version récursive}

\begin{remarque}{Le slicing permet de ne garder qu'un bout de liste}{}

\begin{CodePythontexAlt}[Largeur=1\linewidth,Centre,Lignes=true,PremLigne=1]{}
L = [3, 2, 5, 6, 9, 8, 5]
assert L[:] == L
assert L[2:] == [5, 6, 9, 8, 5]
assert L[:4] == [3, 2, 5, 6]
assert L[2:4] == [5, 6]
            \end{CodePythontexAlt}

\end{remarque}

\begin{exercice}{Écrire la fonction \verb+dichotomie_rec+ avec slicing et la tester}{}


    

\end{exercice}

\section{Diviser pour régner}

\begin{definition}{Diviser pour régner}{}
    Un problème peut se résoudre en employant le paradigme \textbf{Diviser pour Régner} lorsque:

    \begin{itemize}
        \item il est possible de décomposer ce problème en sous problèmes indépendants;
        \item la taille de ces sous problèmes est une fraction du problème final.
    \end{itemize}
    
    La méthode est composée de trois phases:

    \begin{itemize}
        \item la \verb+division+: le problème initial est divisé en sous problème de taille inférieure;
        \item le \verb+règne+: chaque sous-problème est résolu soit récursivement, soit directement si trivial;
        \item la \verb+combinaison+: les solutions des sous problèmes sont combinés pour obtenir la solution du problème initial.
    \end{itemize}

    \includegraphics[width = 0.6\linewidth]{dnc01}
\end{definition}


\begin{exercice}{Algorithme de recherche dans une liste non triée}{}

    \begin{verbatim}
si la liste est vide alors l'élément n'appartient pas à la liste
si la liste ne comporte qu'un seul élément alors l'élément est présent s'il est égal à L[0]
sinon on recherche l'élément dans la liste gauche ou dans la liste droite
    \end{verbatim}

    Écrire cette fonction \verb+recherche(L, val)+ en python
\end{exercice}

\section{Le tri fusion}

Pour trier une paquet de copie dans l'ordre alphabétique, il est possible de scinder le paquet en deux sous-paquets de même taille puis de trier chacun de ces sous-paquets avant de les fusionner pour obtenir le paquet trié.\\

Pour trier ce sous-paquet, on peut utiliser la même technique\dots\\

On peut scinder les sous-paquets jusqu'à avoir une certitude: si le paquet n'a qu'une seule copie alors ce paquet est trié.


\begin{minipage}{0.49\linewidth}
    \begin{verbatim}
    si la liste ne contient qu'un seul élément:
        la liste est triée
    sinon:
        on fusionne les deux sous listes triées

\end{verbatim}
\end{minipage}
\begin{minipage}{0.49\linewidth}
    \includegraphics[width = 1\linewidth]{dnc02}

\end{minipage}



\begin{exercice}{Écrire la fonction \verb+division(L)+ qui renvoie un tuple contenant les deux sous-listes}{}

    \begin{CodePythontexAlt}[Largeur=1\linewidth,Centre,Lignes=true,PremLigne=1]{}
def division(L:list)->tuple:
    
L1 = [1, 4, 6, 15]
L2 = [3, 9, 10]
assert division (L1) == ([1, 4], [6, 15])
assert division (L2) == ([3], [9, 10])
assert division ([1]) == ([] , [1])

                    \end{CodePythontexAlt}
\end{exercice}

\begin{exercice}{Écrire la fonction \verb+fusion(L1, L2)+ qui renvoie une liste fusionnée}{}

\begin{verbatim}
fonction fusion(L1, L2, L)
    si L1 OU L2 est vide:
        on renvoie L1 + L2
    si le premier élément de L1 et plus petit que le premier élément de L2:
        on insère le premier élément de L1 dans la liste fusionnée
    si le premier élément de L2 et plus petit que le premier élément de L1:
        on insère le premier élément de L2 dans la liste fusionnée
    On appelle récursivement fusion avec le nouveau L1, le nouveau L2 et le nouveau L
    
    \end{verbatim}

    \begin{CodePythontexAlt}[Largeur=1\linewidth,Centre,Lignes=true,PremLigne=1]{}
def fusion(L1, L2, L = None):
    if L = None:
        L = []
    return fusion_rec(L1, L2, L)
def fusion_rec(L1, L2, L):
    if L1 == [] or L2 == []:
        return ....
    if ....... < .........:
        .....................
    else:
        .....................
    L = ......................
    return fusion_rec(........, ........, ..........)

assert fusion(L1, L2) == [1, 3, 4, 6, 9, 10, 15]
                    \end{CodePythontexAlt}
\end{exercice}

\begin{exercice}{Ecrire alors la fonction récursive \verb+tri_fusion(L1, L2)+}{}
    \begin{CodePythontexAlt}[Largeur=1\linewidth,Centre,Lignes=true,PremLigne=1]{}
def tri_fusion(L:list)->list:
    
L = [5, 7, 2, 1, 9, 8, 6]
assert tri_fusion(L) == [1, 2, 5, 6, 7, 8, 9]
                            \end{CodePythontexAlt}

\end{exercice}


\begin{propriete}{Complexité de l'algorithme}{}

    \begin{minipage}{0.49\linewidth}
        L'instruction finale 
        
        \verb+fusion(tri_fusion(L1), tri_fusion(L2))+ 
        
        lance deux appels à la fonction \verb+tri_fusion+ (avec certes des données d'entrée deux fois plus petites).

On peut alors montrer que la complexité du tri fusion est en \verb+O(n log n)+
    \end{minipage}\hfill
    \begin{minipage}{0.49\linewidth}
        \includegraphics[width = 0.8\linewidth]{comparaison.png}

    \end{minipage}
    


\end{propriete}

\begin{exercice}{Dessiner de manière schématisé le tri fusion de \verb+[4, 8, 3, 1, 7, 6]+}{}


    \vspace{3cm}
\end{exercice}

\begin{exercice}{Ecrire la trace de cet algorithme lors de l'appel de \verb+tri_fusion([4, 8, 3, 1, 7, 6])+}{}

\begin{tabular}{|p{3.5cm}|c|c|p{1.cm}|p{1.cm}|p{6cm}|p{3cm}|}
    \hline
    Appel &n & IF & L1 & L2 & appel récursif fusion & valeur \\ \hline
    trifusion([4,8,3,1,7,6])&6&F&[4,8,3]&[1,7,6]&trifusion([4,8,3]), trifusion([1,7,6])& \\ \hline
    &&&&&& \\ \hline
    
\end{tabular}
\end{exercice}

    \begin{exercice}{Application à la recherche de maximum dans une liste}{}

        \begin{verbatim}
si la liste ne contient qu'un seul élément:
    c'est le plus grand
on coupe la liste en deux
on récupère le plus grand de la liste de gauche
on récupère le plus grand de la liste de droite
on renvoie le plus grand des deux
        \end{verbatim}
        
        On pourra utiliser la fonction \verb+max+ pour trouver le maximum entre deux valeurs.
        \end{exercice}

    \begin{exercice}{\quad }{}

        Application au calcul de l'exponentiation rapide
        
        Il est facile de vérifier que:
        
        \begin{itemize}
            \item $a^0=1$
            \item si $n$ est pair, $a^n=(a \times a)^{\frac{n}{2}}$
            \item si $n$ est impair: $a^n=a \times (a \times a)^{\frac{n-1}{2}}$
        \end{itemize}
             \end{exercice}


\newpage

\begin{CodePythontexAlt}[Largeur=1\linewidth,Centre,Lignes=true,PremLigne=1]{}

def recherche_dicho_imperative(tab, val):
    ''' Rechercher dichotomique impérative dans liste triée '''
    i_deb = 0
    i_fin = len(tab) - 1
    while i_deb < i_fin:
        i_mil = (i_fin + i_deb)//2
        if tab[i_mil] == val:
            return True
        if tab[i_mil] > val:
            i_fin = i_mil - 1
        else:
            i_deb = i_mil + 1
    return False

def recherche_rec(L, val):
    ''' Recherche dichotomique récursive dans liste triée '''
    if L == []:
        return False
    i_mil = (len(L) - 1)//2
    if L[i_mil] == val:
        return True
    if L[i_mil] > val:
        return recherche_rec(L[:i_mil], val)
    else:
        return recherche_rec(L[i_mil+1:], val)

def recherche_non_triee(L, val):
    ''' Rechercher récursive dans liste non triée '''
    if L == []:
        return False
    if len(L) == 1:
        return L[0] == val
    return recherche_non_triee(L[:len(L)//2], val) or recherche_non_triee(L[len(L)//2:], val)

def division(L):
    ''' Division de la liste en deux sous listes '''
    n = len(L)
    return (L[:n//2], L[n//2:])

def fusion(L1, L2, L = None):
    ''' Fusion de deux listes triées '''
    if L == None:
        L = []
    return fusion_rec(L1, L2, L)
def fusion_rec(L1, L2, L):
    if L1 == [] or L2 == []:
        return L + L1 + L2
    if L1[0] < L2[0]:
        e = L1.pop(0)
    else:
        e = L2.pop(0)
    L = L + [e]
    return fusion_rec(L1, L2, L)

def tri_fusion(L):
    ''' Renvoie la liste L triée par tri fusion '''
    n = len(L)
    if n < 2:
        return L
    else:
        L1, L2 = division(L)
        L1 = tri_fusion(L1)
        L2 = tri_fusion(L2)
        return fusion(L1, L2)
\end{CodePythontexAlt}
    \resetalltcb{}
